\section{在终端里使用 Codex CLI}
\label{sec:ai-cli}

桌面版适合同时看对话、文件和差异；命令行界面（Command Line Interface，CLI）适合已经在终端里构建工程的你，也适合通过 SSH 进入开发机后继续工作。它不是“把网页聊天换成黑底白字”：在获得相应权限后，Codex 可以读取当前工程、修改文件、执行命令，再根据输出继续处理。你要关注的也就不只是最后一句回复，还有它实际操作了哪个目录、改了什么、拿什么证明完成。

\subsubsection{启动后，先确认自己站在哪里}

安装见第~\ref{sec:ai-desktop}~节。第一次使用时，先在普通终端核对程序和登录状态，再进入自己准备的练习工程。通用形式是 \texttt{codex [选项] [初始提示词]}；\texttt{-C 目录} 可以指定工作目录。\texttt{codex login} 会启动登录流程，\texttt{codex login status} 检查是否已有凭据；已登录就不必重复登录。

\Needspace{9\baselineskip}
\begin{tbash}
codex --version
codex login status
codex --help
cd ~/rm-ai-lab
codex
\end{tbash}

这里假设你已把本章练习工程复制到 \path{~/rm-ai-lab}，路径要换成自己的。\texttt{cd} 报错时先处理目录问题，不要站在家目录里继续让代理“随便找项目”。进入界面后，主体显示对话与工具活动，下方是输入框和状态提示。输入 \texttt{/} 会弹出当前版本支持的命令，输入 \texttt{@} 可以查找并引用工程文件。\texttt{@} 帮助模型定位资料，不能代替访问权限，也不等于已经把整个工程永久记住。

用 \texttt{/status} 检查当前模型、目录、权限和上下文使用情况；用 \texttt{/model} 选择可用模型及其支持的思考等级。模型与投入的取舍见第~\ref{sec:ai-model-choice}~节。命令菜单没有某项时，先查看 \texttt{codex --version} 与该版本帮助，不要把一条不支持的命令反复当自然语言发出去。

\begin{example}[第一次对话只读工程]
在练习工程中启动 Codex 后，输入 \texttt{/status}，核对目录；随后在输入框用 \texttt{@} 选中 README，再发送：
\begin{lstlisting}[language={}]
先只读 README 和构建入口，不修改文件、不安装依赖。
说明程序的输入、输出、最短运行命令和你尚未确认的地方。
每项结论给出对应文件位置；找不到时明确说找不到。
\end{lstlisting}
回复后，在 IDE 里打开它引用的文件，核对原文和当前目录，再查看工具记录是否只有读取动作。使用 Git 仓库时，还可以用 \texttt{git status --short} 核对文件状态。完成标准是目录正确、结论有出处、没有修改文件。若它把“README 建议这样运行”说成“已经运行成功”，当场要求更正证据等级。
\end{example}

\begin{exercise}\label{ex:ai-cli-start}
在自己的练习目录中核对 CLI 版本、工作目录与权限，引用一个文件，让 Codex 解释一条真实的构建命令。指出其中项目目标名、输出目录分别是什么，并保留你亲自核对的文件位置；再按下面的方法查一次实际设置来源，区分上下文容量与账号额度。
\end{exercise}

\subsubsection{设置、额度与编辑器上下文怎样查}

\path{~/.codex/config.toml} 保存个人默认模型、沙箱和工具配置；可信项目也可用 \path{.codex/config.toml} 细化本项目设置。命令行参数通常覆盖普通默认值，组织强制策略仍然有效；不可信项目的本地配置层不会照常加载。\texttt{config.toml} 配置客户端如何运行，\texttt{AGENTS.md} 说明项目怎样开发，别把两种文件混成一份。初学时先用选择器与默认设置，需要重复使用某个设置再持久保存，修改现有配置时只改相关条目，别用网上的一份完整配置覆盖自己的环境。

\begin{example}[解释“我明明改过设置，怎么没生效”]
空闲时输入 \texttt{/debug-config}，检查配置层、当前工作目录和生效策略，再用 \texttt{/status} 对照实际模型与权限。如果项目配置覆盖了个人默认值，或组织策略禁止某项权限，就已找到调查方向；不要靠反复重启猜原因。输出可能含本机路径与服务信息，求助时先脱敏。

随后用 \texttt{/usage} 查看账号允许显示的用量与限额信息。\textbf{账号额度与上下文剩余容量不同}：压缩上下文能腾出工作空间，不会充值账号，也不会清除已经产生的 API 费用。原始 API Key 登录时，账号用量菜单可能不提供同类统计，应查看对应 Platform 用量与计费页面。

使用已连接的受支持 IDE 时，\texttt{/ide} 可把打开文件、选区等编辑器上下文加入下一条请求；发送前仍要指出要处理的代码与目标。普通 SSH 终端没有相应集成时，直接用 \texttt{@} 引用文件并提供位置，不要假定代理看见了你另一台机器的编辑器。
\end{example}

按需用 \texttt{/export} 导出聊天记录，或用 \texttt{/recap} 生成回顾；导出要核对是否带有私有日志，回顾也不等于压缩上下文。更进阶的\textbf{生命周期钩子（Hooks）}可以在工具调用前后、轮次结束等事件触发脚本或 MCP 工具，适合已经稳定的检查流程。它与按任务选用的 Skill 不同，可能自动执行；先用 \texttt{/hooks} 查看来源和动作，新改的非托管钩子需要信任审查。不要把还没手动跑通的命令塞进每一轮自动执行。设置与钩子的具体配置见\href{https://learn.chatgpt.com/docs/config-file/config-basic}{Config basics}和\href{https://learn.chatgpt.com/docs/hooks}{Hooks 官方说明}；本章先掌握查看与判断，编写复杂钩子留待真实需要时。

\subsubsection{权限、搜索与运行中的干预}

在本例 Codex 中，先分清两层：\textbf{沙箱（Sandbox）}规定工具技术上能访问什么；\textbf{审批策略（Approval Policy）}规定哪些动作要先问你。项目内自动修改不代表整个电脑都交给它；“不要改文件”的提示词也不能替代只读沙箱。刚接触陌生工程，先用 \texttt{/permissions} 选择只读预设；开始已确认的小任务，再使用允许工程内写入的配置。换用其他工具时，核对它实际提供的访问限制与审批机制，不能假定也有相同预设。

下面是一种明确指定边界的启动方式：
\begin{tbash}
codex --sandbox workspace-write --ask-for-approval on-request
\end{tbash}
\texttt{workspace-write} 允许在配置的工作区内写入，\texttt{on-request} 允许代理在需要时提出审批；组织管理策略还可能进一步限制它。遇到审批框，读清命令、工作目录、目标文件和联网目的，再决定是否批准。安装依赖、修改系统设置和访问外部服务不能靠“看起来是 AI 提出的”就一路同意。

查询当前 API 或软件安装方法时，要求它读取官方文档并给出链接。CLI 的 \texttt{--search} 用于开启实时网络搜索，默认搜索模式可能使用缓存；\textbf{有搜索能力不等于终端命令已经获得任意联网权限}。如果代理说资料已过时，先核对来源日期、依赖版本与实际接口，而不是立即升级整个系统。

\begin{example}[只给当前任务需要的权限]
先在只读模式要求 Codex：“查出测试入口，说明要运行哪些命令、它们会写入哪里；暂不运行。”确认它只需在工程构建目录写文件后，切到工程写入模式，再要求执行这些检查。若出现外部目录或网络审批，逐项核对理由；无需批准与任务无关的动作。预期结果是得到实际命令、退出状态和关键输出，而非一句“应该没问题”。
\end{example}

运行中发现方向不对，可以直接补充明确的限制。当前 CLI 文档规定：任务执行时按 \texttt{Enter} 发送的新指令会进入当前轮，按 \texttt{Tab} 则排队到下一轮；以你版本的输入框提示为准。需要停下来时，使用界面显示的中断操作；若启用了目标模式，先暂停目标，防止后续继续推进。\texttt{/ps} 查看当前会话启动的后台终端，\texttt{/stop} 停止这些后台终端；关闭对话不能当成“所有服务器和训练进程都已消失”的证据。

\begin{note}
停止生成、停止终端、暂停目标、撤销改动是四件事。停止之后，已经写入的文件、已经发送的请求仍然存在。不要为了省一次审批把“绕过全部审批与沙箱”作为日常启动参数；需要额外目录时，先核对 \texttt{--add-dir 目录} 是否已能满足任务。
\end{note}

\begin{exercise}\label{ex:ai-cli-permissions}
Codex 为了运行项目测试，提出“安装一个系统软件包”和“删除整个项目重新下载”两个动作。分别写出批准前需要核对的信息。再说明停止回复后，为什么仍要检查后台终端和已经产生的文件改动？
\end{exercise}

\section{按任务使用计划、目标与会话}
\label{sec:ai-workflow}

碰到要分几步才能完成的任务时，先划清当前这一轮要解决的问题。你可能同时想修串口接收、换一个滤波器、顺手重构整个工程。如果都放进一个无限延长的聊天里，需求会互相覆盖，改动也难以验收。我们先决定这一次究竟交付什么，再选择控制它的工具。

先规划、明确验收终点、保存可撤销的版本、隔离并行改动，都是可以迁移的方法。下面的 Plan、Goal 与会话入口以 ChatGPT 桌面版及 Codex CLI 为例；其他工具未必有同名模式。没有持久目标功能时，可以用任务卡与分阶段检查推进，再由你决定是否继续；长任务续跑来自 Harness 的实现，并非 LLM 自带的保证。

\subsubsection{计划模式先定路线，目标模式持续检查终点}

\textbf{计划模式（Plan Mode）}用于动手前调查现状、问清约束并给出实施路线。涉及多个模块、存在几条可选方案、迁移可能改变接口，或你自己还说不清需求时，先用 \texttt{/plan}。一个好的计划应写明：现状证据、改动范围、分步实施、每步如何验证、失败时怎样退回。计划中出现“优化架构”“完善系统”，却没有文件、接口或检查办法，还需要继续收窄。

桌面版可通过 \texttt{/plan} 切换；CLI 在空闲时也可输入 \texttt{/plan 任务描述} 进入计划模式。完成讨论后，使用界面提供的实施入口或模式切换，确认已回到执行模式，再授权实施。支持相应快捷键的版本也可使用 \texttt{Shift+Tab}；看当前模式提示，别只凭键按过了就认定切换成功。\textbf{计划模式不是“输出一个步骤列表”的同义词}，普通执行中的进度清单也不代表代理处于计划模式。

\begin{example}[先规划一次数值换算的修复]
在练习工程中进入计划模式，发送：
\begin{lstlisting}[language={}]
先读取毫米换算函数、接口、task.md 和测试，不修改文件。
说明输入输出类型、整数除法与精度丢失发生的位置。
核对 0、1000、1500、-500 mm 的期望结果，负位移是合法输入。
给出最小修改计划、回归检查和撤回办法，暂不改文件。
\end{lstlisting}
你应该先得到真实现状：现有函数接收整数毫米、返回浮点米，但除法已在返回前丢失了小数。核对它是否读了测试、是否承认负位移合法，确认计划没有带上无关输入解析或重构，再切回执行模式要求完成这一项。若需求仍没定，继续讨论即可。
\end{example}

\textbf{目标模式（Goal Mode）}把可验证的终点挂在当前对话上，让代理在多轮工作中持续推进。适合“复现偶发错误，修复后让相关测试稳定通过”这种下一步取决于上一轮证据的任务；改一句注释通常用普通请求就够了。它的关键是\textbf{结果、约束、验证、停止条件}，不是“永远干下去”。

\Needspace{12\baselineskip}
\begin{example}[给多轮修复设置有边界的目标]
在练习工程中已确认换算契约后，可以输入：
\begin{lstlisting}[language={}]
/goal 修复本练习工程毫米换算米时的小数截断，保持公开接口，
保留负位移语义。让原有 0、1000、1500、-500 mm 用例通过，
另核对 1、999、-1500 mm；只改相关实现，不删除测试或放宽精度。
按 README 重新构建并运行同一组测试，记录实际命令和结果。
如果缺少依赖或需要改变契约，先报告证据与所需决定。
不得把未运行的检查写成通过。
\end{lstlisting}
CLI 用 \texttt{/goal} 查看目标，\texttt{/goal edit} 修改文字，用 \texttt{/goal pause}、\texttt{/goal resume}、\texttt{/goal clear} 暂停、恢复或清除目标。桌面版的目标进度行位于输入框上方，可在那里暂停、恢复、编辑或清除。练习时至少演示一次暂停，再核对目标状态，避免把“我看到它暂时没输出”误当成暂停。
\end{example}

目标跟随当前对话，不会自动变成所有项目的永久规则。它只在允许的状态与资源预算内推进；需要你决定、触及权限边界、被中断或预算耗尽时都可能停止。\textbf{预算用完和目标完成是不同状态}。如果客户端提供 token 预算设置，应在启动前设置可接受的预算，并要求停下时交代已完成内容、证据与剩余工作。本次核对的用户文档没有给出统一的预算命令或调整入口；是否能在运行后调整，以客户端实际控件为准，不要照抄网上未经确认的 \texttt{/goal budget} 写法。目标也不会解除沙箱、接上不存在的设备，或保证关机断网后照常工作。

\begin{conclusion}
普通请求适合明确的一次任务；计划模式适合先把路线与边界议定；目标模式适合终点清楚、路径需要反复试验的工作。可以先计划，再执行，再按需要设目标。无论用哪个入口，交付标准都应由证据决定。
\end{conclusion}

\begin{exercise}\label{ex:ai-workflow-plan-goal}
为“改一处变量名”“迁移相机 SDK”“定位只在板端偶发的崩溃”分别选择普通请求、计划模式或目标模式，并说明何时需要组合使用。再给第三项写一个目标：没有板端日志或设备时，应当报告什么，不能声称什么？
\end{exercise}

\subsubsection{恢复与分叉对话}
\label{sec:ai-dialog-history}

关闭 CLI 后继续昨天的工作，使用 \texttt{codex resume} 打开历史选择器，或在工程目录执行 \texttt{codex resume --last}。\texttt{/resume} 也能在交互界面恢复已有会话。恢复的是对话；仓库文件可能已被别人修改，所以回来后要重新核对当前源码与工作目录。

\textbf{分叉（Fork）}从已有历史创建另一条对话，适合保留一条路线再探索另一条。桌面版或 CLI 输入 \texttt{/fork}，按界面选择当前目录或新 worktree；退出 CLI 后也可用 \texttt{codex fork} 选择历史。\texttt{/new} 则从新的对话上下文开始，适合已经换了任务。短暂询问术语或让另一路复核一个想法，可以用支持版本中的 \texttt{/side} 侧聊；侧聊结果要影响主任务时，明确把决定带回主对话。

分叉对话和 Git 分支是两条轴：前者分开消息历史，后者记录代码版本。\textbf{两个分叉如果仍选同一个目录，它们看到并修改的仍是同一套文件。}需要让两路各自修改时，再查第~\ref{sec:ai-project-workflow}~节的工作树练习。

所谓“回滚对话”，通常是回到某个较早的消息位置，修改当时的指令并从那里重走，或移除后面的历史。适合你一开始就把需求讲错、后续讨论已沿错误前提展开的情况。当前 CLI 提供一个具体入口：\textbf{输入框为空时连续按两次 \texttt{Esc}}，进入编辑上一条用户消息并从该位置分叉的流程。核对被编辑的原消息，改正前提后再发送；原路线可用历史选择器保留查阅。若只是补充一个条件，直接追加纠正通常更方便。

桌面版用 \texttt{/fork} 保留当前路线；某些版本还在消息菜单提供编辑、从此处分叉或历史回退操作，具体以当前菜单与确认框为准。不要把旧教程中的按钮名当成每个版本都有的功能，更不要混用面向程序开发者的内部接口来“修复”对话。


\begin{note}
继续、分叉或回退对话之后，都要核对实际文件与后台任务。需要撤销文件改动时，先看第~\ref{sec:ai-project-workflow}~节的版本检查与恢复练习。
\end{note}

\begin{example}[无副作用地试一次“回到上条指令”]
在一个没有运行任务的练习对话中，先要求：“只给出两种队列策略的文字比较，不改文件。”得到回复后，让输入框保持空白，连续按两次 \texttt{Esc}，把原消息改成“只比较容量为一的最新帧队列和先进先出队列，并说明丢帧语义”。发送后检查两条路线的历史和工具记录，再打开练习文件核对内容。这个例子练的是改前提和分叉。
\end{example}

\subsubsection{多代理并行：把边界写到任务里}

子代理（Subagent）可以把可独立完成的工作分出去，例如一路查接口，一路读日志，一路审查测试缺口，由主代理汇总。当前 Codex 客户端可以直接用自然语言要求并行；CLI 的 \texttt{/subagents} 可以查看和切换子代理对话。子代理各自消耗上下文与模型资源，开得更多不一定更快：工作本身难以分开时，协调与冲突会抵消收益。

先从只读任务开始最容易观察效果。写入任务则要明确每个代理的文件范围、输入与交付物；有依赖的步骤按先后执行。\textbf{子代理不等于自动拥有独立 worktree}，也不能因为名叫“审查者”就认为权限一定是只读。本例的子代理通常继承父任务权限，其他系统可能分别配置；“只读检查”的职责仍要明确交代，技术隔离则通过相应执行环境完成。

\begin{example}[并行找证据，统一决定修改]
在一个已保存当前工作的工程中发送：
\begin{lstlisting}[language={}]
请用两个子代理并行只读检查，不修改文件、不访问真实设备。
A：检查输入解析和调用方，报告接口约定与边界风险。
B：检查现有测试，报告已覆盖与未覆盖的边界。
每一路返回文件位置、可复现输入和不确定项。
主代理等待两路结果后，合并重复问题，给出最小修复建议；
没有证据的问题标成待验证，不直接改代码。
\end{lstlisting}
查看两路是否按职责返回，主代理是否保留了证据出处。决定实现后，再给一个明确任务去修改与验证。这个过程覆盖并行分工、结果汇总和写入边界，不需要让三个代理同时抢一个文件。
\end{example}

\begin{exercise}\label{ex:ai-workflow-parallel}
为“更换串口解析模块”设计三个并行只读子任务，写明每个任务的输入、输出和禁止修改范围。再指出哪些实现步骤必须等其他结果完成，不能直接并行。
\end{exercise}

\begin{note}
本节的 CLI 操作按 2026 年 10 月 31 日官方文档核对，本机只读检查版本为 \texttt{0.166.0}。命令清单与快捷键见\href{https://learn.chatgpt.com/docs/developer-commands}{官方 Developer commands}，桌面模式入口见\href{https://learn.chatgpt.com/docs/reference/slash-commands}{桌面斜杠命令}；计划模式见\href{https://learn.chatgpt.com/guides/best-practices}{Best practices}。目标的生命周期与停止条件见\href{https://developers.openai.com/cookbook/examples/codex/using_goals_in_codex}{Goals 使用指南}及\href{https://learn.chatgpt.com/docs/long-running-work}{长任务说明}。工作树和子代理分别见\href{https://learn.chatgpt.com/docs/environments/git-worktrees}{Worktrees}、\href{https://learn.chatgpt.com/docs/agent-configuration/subagents}{Subagents}；会话回退的文件边界见\href{https://learn.chatgpt.com/docs/app-server\#revert-thread-history}{历史回退说明}。这些链接用于随版本复核，不要求新手实现内部协议。
\end{note}
